\begin{frame}{자주 묻는 질문}
  \small
  \begin{itemize}
    \item \textbf{하이퍼파라미터 선택 근거까지 물어봐야 하나?}
      \textbf{그렇다} — ``적당히'' 고른 값($\gamma, \lambda$,
      $\epsilon$, 보상 스케일)이 결과에 얼마나 큰 영향을 주는지 확인
      했는지가 \textbf{방법론 적절성} 평가의 핵심.
    \item \textbf{모방학습/MCTS 팀은 체크리스트를 어떻게?} ``학습
      곡선'' 대신 다른 증거로 — 모방학습: Ch12.1 \textbf{복합
      오차}(전문가 분포 밖 성능 저하) 관찰됐는가 / MCTS: Ch15.2처럼
      \textbf{시뮬레이션 횟수 증가} 시 선택되는 수가 더 정확해지는가.
    \item \textbf{문제가 없는 팀에 3점 질문을 어디서?} ``문제가
      \emph{보이지 않는} 것''과 ``\emph{확인되지 않은} 것''은 다르다 —
      3점 질문은 항상 \textbf{확인되지 않은 영역}에서 나옴(시드 5개
      확장, $\lambda$ 민감성, 더 긴 학습 스텝, 다른 평가 프로토콜).
      ``이 실험의 다음 단계로 무엇이 가장 정보가 많을까요?''가 3점이
      될 수 있다 — 발표팀이 답할 수 있다면 발표는 \textbf{더 단단해}
      짐(리뷰어가 목표 달성).
    \item \textbf{``더 오래 돌려보세요''를 쉽게 요구할 수 없는데?}
      리뷰어가 직접 300만 스텝을 돌릴 필요는 \textbf{없다} — 3점
      질문은 \emph{발표팀}이 돌려야 할 실험을 \emph{지적}하는 것.
      ``12만 스텝으로 포화가 안 보이는데, 이 과제는 완전 수렴에 수십만
      스텝이 필요하다고(11.2) 했죠? 그걸 \emph{인지하고} ``아직
      안 끝났다''고 보고했나요?'' — ``몰랐습니다''면 \textbf{수렴
      여부 미인식} 자체가 ``결과의 정직성'' 감점 요인.
  \end{itemize}
\end{frame}
